\begin{trapbox}

	\begin{description}[style=nextline, leftmargin=2.8em, labelsep=0.5em, itemsep=0.35em]
		\item[\textbf{\textcolor{CTrap}{易错点一（忽略端点）：}}]
			学生只关注 $f'(x)=0$ 的点，忘记计算区间端点 $a,b$ 的函数值，导致最值遗漏。
		\item[\textbf{\textcolor{CTrap}{正确理解（比较所有可疑点）：}}]
			最值候选点必须包括端点、导数为零的点和不可导点，不能漏掉任何一个。
		\item[\textbf{\textcolor{CTrap}{错因分析（思维定势）：}}]
			常误以为最值只出现在极值点，忽略端点也可能是最值点。
	\end{description}

	\medskip
	\begin{description}[style=nextline, leftmargin=2.8em, labelsep=0.5em, itemsep=0.35em]
		\item[\textbf{\textcolor{CTrap}{易错点二（忽视不可导点）：}}]
			求最值时只解 $f'(x)=0$，不考虑导数不存在的点，例如绝对值函数的转折点。
		\item[\textbf{\textcolor{CTrap}{正确理解（检查不可导性）：}}]
			必须同时检查 $f'(x)$ 为零和 $f'(x)$ 不存在的点，两者都是可疑点。
		\item[\textbf{\textcolor{CTrap}{错因分析（条件遗漏）：}}]
			对“可导（允许有限个不可导点）”条件理解不深，忽略不可导点的存在。
	\end{description}

	\medskip
	\begin{description}[style=nextline, leftmargin=2.8em, labelsep=0.5em, itemsep=0.35em]
		\item[\textbf{\textcolor{CTrap}{易错点三（单调性误判）：}}]
			不验证单调性，随意假定函数单调递增或递减，从而错误地只比较端点函数值，忽略内部可疑点（驻点、不可导点）。
		\item[\textbf{\textcolor{CTrap}{正确理解（结合单调判断）：}}]
			单调递增时最小值在左端点，最大值在右端点；单调递减时相反。但若非单调，必须全面比较。
		\item[\textbf{\textcolor{CTrap}{错因分析（思维片面）：}}]
			只考虑了一种单调情形，未进行整体分类讨论。
	\end{description}

\end{trapbox}
